%% =====================================================================================
%% Revision kit for the IET EPA submission: every text change as ready-to-apply LaTeX.
%%
%% Each block states where it goes and what it replaces. Equation references are written
%% as literal numbers, e.g. (5); when applied to the manuscript source they are mapped to
%% its existing \label/\eqref names. Items marked [AUTHORS] need information only the
%% authors have. See review notes in analysis/08_iet_epa_presubmission_review.md.
%% =====================================================================================


%% -------------------------------------------------------------------------------------
%% A6. ABSTRACT  - replaces the whole abstract (~200 words)
%% -------------------------------------------------------------------------------------
Reliable diagnosis of rotor winding inter-turn short circuits (ITSCs) in synchronous
generators is challenging under operating-point variation, because excitation and reactive
power strongly affect stray magnetic field measurements. This paper presents a
leakage-safe cross-operating-point evaluation of rotor winding short-circuit severity
diagnosis using two-channel stray-field signals from a grid-connected salient-pole
synchronous generator at seven excitation and power-factor conditions. Physics-informed
feature representations are evaluated with classical learners and raw-signal 1D-CNN models
under a leave-one-operating-point-out (LOOPO) protocol. Fundamental-normalized descriptors
are far more robust than absolute magnitude-dependent features (macro-F1 of 0.9455 versus
0.8026 with the same classifier); the normalization cancels excitation dependence almost
completely in over-excited operation but only partially in under-excited operation,
identifying the under-excited regime as the diagnostic boundary. ExtraTrees with all
numeric descriptors attains a macro-F1 of 0.9951 (0.49\% file-level error), while SVM-RBF,
logistic regression and a nested severity-calibrated 1D-CNN each attain 0.9935 (0.66\%),
and raw 1D-CNN majority voting attains 0.9737 (2.64\%). At the most under-excited
condition, where the raw CNN misreads 20\% faults as healthy, nested ordinal severity
calibration corrects every such error without degrading any correct prediction. Robust
stray-field rotor winding diagnosis therefore depends on physically meaningful
normalization and leakage-safe evaluation at least as much as on model complexity.


%% -------------------------------------------------------------------------------------
%% B1. SECTION 2 - replaces the sentence that begins
%%     "The model assumes linear magnetics and a fixed defect location and neglects local
%%      saturation and AVR dynamics, ..." (through "... is designed to test.")
%% -------------------------------------------------------------------------------------
The model assumes linear magnetics and a fixed defect location, and neglects local
saturation, AVR dynamics and the contribution of the grid-imposed terminal voltage and
armature current to the fundamental. The first two perturb the constants $K_k$ but not the
leading-order structure of (5). The third does not: in under-excited operation it limits
the cancellation of $I_\mathrm{f}$ in (5), as shown in Section~5.1. Because the public
dataset does not report the field current at each operating point, the excitation ordering
used here follows from standard phasor relations for a grid-connected synchronous generator
at constant apparent power rather than from direct measurement. Whether the first-order
reduction of excitation dependence is sufficient under real operating-point variation is
exactly what the leakage-safe benchmark of Section~4 is designed to test.


%% -------------------------------------------------------------------------------------
%% B2. SECTION 5.1 - insert as a new paragraph directly after
%%     "However, the collapse is not complete: visible residual spread remains between
%%      under-excited and over-excited conditions."
%%     The 2.5x figure is read from Fig. 5(c) (P18 ~0.29 vs. P19-P21 ~0.73 at gamma = 1);
%%     [AUTHORS] confirm against the plotted values if convenient.
%% -------------------------------------------------------------------------------------
The structure of this residual is informative. The three over-excited operating points
(P19--P21) collapse almost onto a single curve, whereas the under-excited points fan out
monotonically with decreasing excitation from P15 to P18, reaching a spread of roughly
$2.5\times$ at $\gamma = 1$; a single linear fit pooled over all operating points gives
$R^2 = 0.811$. This pattern follows if the fundamental contains a component that does not
scale with field current, which (4) omits. Because the generator is grid-connected at
constant apparent power, its terminal voltage and armature-current magnitude are imposed
externally, so part of the 60-Hz stray field is fixed by the operating condition rather
than by $I_\mathrm{f}$. Writing $|V(f_\mathrm{e})| \approx K_a + K_\mathrm{e} I_\mathrm{f}$
gives
\[
  R(kf_\mathrm{m}; f_\mathrm{e}) \approx
  \frac{\gamma K_k I_\mathrm{f}}{K_a + K_\mathrm{e} I_\mathrm{f}} .
\]
At high excitation the rotor term dominates and $R$ tends to the excitation-independent
value $(K_k/K_\mathrm{e})\gamma$ predicted by (5), whereas at low excitation $R$ falls with
$I_\mathrm{f}$. Normalization therefore removes the excitation dependence in over-excited
operation but only partially in under-excited operation, which is precisely the regime of
the boundary failure analysed in Section~5.4.


%% -------------------------------------------------------------------------------------
%% A5. SECTION 5.2 - insert directly before the sentence
%%     "Moreover, macro-F1 depends on the class-wise distribution of errors ..."
%% -------------------------------------------------------------------------------------
For this reason, Table~4 and Figure~6 report error rates, which are comparable across all
rows.


%% -------------------------------------------------------------------------------------
%% C. SECTION 5.4 - replaces the two sentences
%%    "The selected upper threshold was stable across all seven outer LOOPO folds, ...
%%     ... than by fragile threshold tuning."
%% -------------------------------------------------------------------------------------
The P18 corrections return 20\% faults from the healthy class to the 20\% class, a decision
governed by $t_{01}$ and $t_{12}$. Both thresholds were fixed a priori at the physical
midpoints between severity levels and were never tuned, so the correction cannot be
attributed to threshold selection. The tuned threshold $t_{23}$ affects only the 50\%/100\%
boundary; it was stable at 0.675 across all seven outer LOOPO folds, and P18 performance
remained perfect for $t_{23}$ between 0.55 and 0.80. The correction is therefore better
explained by ordinal severity information than by threshold tuning.


%% -------------------------------------------------------------------------------------
%% D2. SECTION 3 - append to the paragraph that introduces P15-P21
%% -------------------------------------------------------------------------------------
The public dataset comprises 607 two-channel voltage-vector files, of which 95, 84, 80, 90,
90, 84 and 84 correspond to P15--P21, respectively; all of them are used in this study.


%% -------------------------------------------------------------------------------------
%% A3 + A5. TABLE 4 - full-width (table*) so it no longer collides with Table 3 on
%%          page 9, with an error-rate column. Match the manuscript's existing table
%%          markup (the Wiley NJD template's tablenotes environment) when applying.
%% -------------------------------------------------------------------------------------
\begin{table*}[t]
\caption{Strong classical and raw deep benchmark under LOOPO.}
\label{tab:benchmark}
\centering
\begin{tabular}{llrccccr}
\hline
Model & Representation & $n$ & F1 & Bal. & MAE & Err. & Err. rate \\
\hline
ET      & All numeric       &  607 & \textbf{0.9951} & \textbf{0.9953} & \textbf{0.0010} &  3 & \textbf{0.49\%} \\
SVM     & Norm. spectral    &  607 & 0.9935 & 0.9937 & 0.0015 &  4 & 0.66\% \\
LR-L2   & Band-energy       &  607 & 0.9935 & 0.9936 & 0.0013 &  4 & 0.66\% \\
SC-CNN  & Raw + calibration & 1821 & 0.9935 & 0.9936 & 0.0013 & 12 & 0.66\% \\
Raw-CNN & Raw voting        & 1821 & 0.9737 & 0.9733 & 0.0053 & 48 & 2.64\% \\
RF      & Cross-channel     &  607 & 0.9632 & 0.9619 & 0.0157 & 22 & 3.62\% \\
\hline
\end{tabular}
\par\smallskip
\footnotesize ET: ExtraTrees; LR-L2: logistic regression with $\ell_2$ regularization;
SC-CNN: nested severity-calibrated 1D-CNN. $n$: file-level predictions (classical rows
cover the 607 files once; CNN rows pool three seeds, $3\times607=1821$). F1: macro-F1;
Bal.: balanced accuracy; MAE: severity mean absolute error; Err.: misclassified
predictions. Error rates, unlike error counts, are comparable across rows.
\end{table*}


%% -------------------------------------------------------------------------------------
%% A2. FIGURE CAPTIONS
%%     [AUTHORS] The attribution assumes Figs. 1-3 are reproduced from the data report
%%     [14]. If you drew any of them yourselves, drop its last sentence.
%% -------------------------------------------------------------------------------------

% Figure 1 (workbench)
\caption{Experimental workbench used to acquire the public dataset: (1) 10~kW DC prime
mover; (2) 8-pole salient-pole synchronous generator analysed in this study; (3) 2-pole
cylindrical-rotor synchronous generator, not used here; (4) electromagnetic clutch;
(5) mechanical unbalance/flange systems; (6) position transducers. Only the 8-pole
salient-pole generator was used for the rotor-winding ITSC measurements analysed here.
Reproduced from [14] under the CC~BY~4.0 licence.}

% Figure 2 (tap configuration)
\caption{Rotor winding configuration showing the taps on a single pole used to impose 20\%,
50\% and 100\% short-circuit severities (94, 235 and 470 of the pole's 470 turns), selected
through terminals J1, J2 and J3; terminal J corresponds to the healthy winding. Reproduced
from [14] under the CC~BY~4.0 licence.}

% Figure 3 (sensor)
\caption{(Left) Induction sensor (search coil); (right) 90$^\circ$ spatial placement of CH1
and CH2 on the generator's exterior frame. Reproduced from [14] under the CC~BY~4.0
licence.}

% Figure 6 (use revision/figures/fig6_macro_f1.pdf)
\caption{Macro-F1 of strong classical learners and of raw-signal and severity-calibrated
CNN decisions under the LOOPO protocol. Labels give macro-F1 and the file-level error rate.
Classical results cover the 607 files; CNN results pool three seeds (1821 file-level
predictions), so error rates rather than error counts are comparable across bars. SC-CNN
denotes the nested severity-calibrated 1D-CNN.}

% Figure 7 (P18 confusion matrices)
\caption{P18 boundary-condition confusion matrices before and after nested severity
calibration. Counts are pooled over three seeds, so each row total is three times the
per-seed number of files (P18 contains 30 healthy files and 20 files at each fault
severity). The raw 1D-CNN mainly misclassifies 20\% faults as healthy, whereas nested
severity calibration recovers all P18 errors.}


%% -------------------------------------------------------------------------------------
%% A1. ACKNOWLEDGMENTS - new section, placed before "Conflicts of Interest"
%%     IET EPA asks for all funding sources to be listed here.
%%     [AUTHORS] add the grant number and confirm the funder name.
%% -------------------------------------------------------------------------------------
\section*{Acknowledgments}
This work was supported by the Research Council of Finland (formerly the Academy of
Finland) under the Centre of Excellence Programme for the project High-Speed
Electromechanical Energy Conversion Systems [grant number]. The authors thank the authors
of [14] for making the stray magnetic field dataset openly available.


%% -------------------------------------------------------------------------------------
%% LAYOUT (applied in the source; no text change)
%%   - Figure 5: make it full width (figure*) so its legends and axes are legible.
%%   - Figure 6: replace the image with revision/figures/fig6_macro_f1.pdf.
%%   - References: switch to revision/references.bib. The double commas (",, and") are
%%     produced by the bibliography style, not the .bib entries, so the style is fixed too.
%% -------------------------------------------------------------------------------------
